\section{CRF 对整条标签序列做条件归一化}
\label{sec:13-Machine-Learning/07-Graphical-and-Sequential-Models/04}

你在做一个简历解析器，要把每个词标成 B-ORG、I-ORG、B-TITLE、O 之类的标签。训练完一看，逐词准确率 \(96\%\)，指标很体面。可打开输出，问题一眼就能看见：O 后面直接跟着 I-ORG，一个实体被切成两段，另一处又把两家公司粘成了一个。逐词准确率高，整句全对的比例却低得多——二十个词都独立判对的概率是 \(0.96^{20}\approx 0.44\)。分类器把每个位置当成独立决策，而标签之间显然有约束。

\subsection{生成式模型要为输入建模，判别式模型不必}

隐 Markov 模型早就能表达标签之间的依赖，它写下的是联合分布

\[
p(x_{1:T},y_{1:T})
=p(y_1)\prod_{t=2}^{T}p(y_t\given y_{t-1})
\prod_{t=1}^{T}p(x_t\given y_t),
\]

细节见\hyperref[sec:13-Machine-Learning/06-Probabilistic-Models/05]{``HMM 把隐变量排列成时间链''}。它能对整条路径解码，但也必须为输入建模：那个 \(p(x_t\given y_t)\) 要求给定标签之后，各位置的观测彼此独立。

在文本任务里这一条几乎肯定是假的。你想用的特征包括词形、前后缀、是否大写、词典命中、上下文窗口里的词、预训练向量——它们高度重叠、彼此强相关。给定标签后它们条件独立的说法，是为了让联合分布写得出来才成立的，不是任务相信的结构。而解析器根本不需要生成简历文本，它只需要在文本已经摆在面前时给出标签。

既然如此，就只对需要的那部分建模。

\subsection{全局归一化把整条路径绑在一起}

线性链条件随机场直接写条件分布：

\[
p_\theta(y_{1:T}\given x_{1:T})
=\frac{\exp s_\theta(x,y)}{Z_\theta(x)},
\qquad
s_\theta(x,y)=\sum_{t=1}^{T}\theta^{\trans}f_t(y_{t-1},y_t,x,t),
\]

而 \(Z_\theta(x)=\sum_{y'}\exp s_\theta(x,y')\) 对所有完整标签序列求和。特征 \(f_t\) 可以随意读取整段输入、当前位置和相邻标签，不必声称这些输入分量之间有任何独立性——因为 \(x\) 从头到尾都站在条件线右边，模型不对它的分布负责。这是与 softmax 分类器同一族的对数线性模型（见\hyperref[sec:13-Machine-Learning/02-Linear-Models/03]{``从二分类到 Softmax 分类''}），只是把``类别''换成了``整条路径''。

失去的东西也在这里。CRF 给不出 \(p(x)\)，不能生成文本，也无法比较两段输入本身谁更常见；缺失输入、半监督利用未标注数据这些事，它都不如生成式模型顺手。

归一化常数依赖于 \(x\)，并且是对全体路径求和。这一点把局部选择绑成了整体：抬高某条路径的得分，同一输入下别的路径就相对被压低。标签数为 \(K\)、长度为 \(T\) 时路径有 \(K^{T}\) 条，直接枚举当然不可行。

\subsection{前向递推把路径求和压成边上的遍历}

把局部势记作 \(\psi_t(i,j;x)=\exp\bigl(\theta^{\trans}f_t(i,j,x,t)\bigr)\)，定义走到位置 \(t\)、末标签为 \(j\) 的全部前缀的总权重

\[
\alpha_t(j)=\sum_{y_{1:t-1}}\prod_{r=1}^{t}\psi_r(y_{r-1},y_r;x).
\]

它满足

\[
\alpha_t(j)=\sum_{i}\alpha_{t-1}(i)\,\psi_t(i,j;x),
\qquad
Z_\theta(x)=\sum_j\alpha_T(j).
\]

递推之所以成立，靠的是链结构给出的条件独立：给定 \(y_{t-1}\) 之后，前缀怎么走完全不影响后缀的权重，于是所有以 \(i\) 结尾的前缀可以合并成一个数。

\begin{center}
\begin{tikzpicture}[line width=0.7pt]
% 转移边
\foreach \t in {0,1,2,3}{
  \foreach \i in {0,1,2}{
    \foreach \j in {0,1,2}{
      \draw[black!22] (1.6*\t+0.22,\i) -- (1.6*\t+1.38,\j);
    }
  }
}
% 高亮一条路径
\draw[line width=1.4pt] (0.22,2) -- (1.38,0);
\draw[line width=1.4pt] (1.82,0) -- (2.98,0);
\draw[line width=1.4pt] (3.42,0) -- (4.58,1);
\draw[line width=1.4pt] (5.02,1) -- (6.18,2);
% 节点
\foreach \t in {0,1,2,3,4}{
  \foreach \i in {0,1,2}{
    \fill[white] (1.6*\t,\i) circle (0.22);
    \draw (1.6*\t,\i) circle (0.22);
  }
}
% 行列标注
\foreach \i/\lab in {2/{B-ORG},1/{I-ORG},0/{O}}{
  \node[anchor=east,font=\small] at (-0.45,\i) {\lab};
}
\foreach \t/\lab in {0/1,1/2,2/3,3/4,4/5}{
  \node[font=\small] at (1.6*\t,-0.65) {\(t=\lab\)};
}
\node[anchor=west,font=\small] at (6.9,1.6) {路径共 \(K^{T}=243\) 条};
\node[anchor=west,font=\small] at (6.9,0.9) {边只有 \((T-1)K^{2}=36\) 条};
\end{tikzpicture}
\end{center}

\emph{前向递推沿边走一遍就够，不必逐条枚举路径。}

图里 \(K=3\)、\(T=5\)，路径二百四十三条，边三十六条。递推的每一步只沿着一列到下一列的 \(K^2\) 条边做加权求和，总开销 \(O(TK^2)\)——这正是上一节所说的``链的树宽为一''在具体算法里的样子，Kalman filter 的预测-更新与它是同一套局部消元。后向递推对称地定义，前向量与后向量相乘就得到某个位置或某对相邻标签的边缘概率。

实现时递推要放在对数域里用 \(\logsumexp\) 做，否则几十个小于一的势连乘之后会直接下溢成零。

\subsection{求和与取最大是两种不同的动态规划}

同一张格子上跑两种递推，用途完全不同，混起来是常见错误。前向算法用求和合并路径，得到的是归一化常数与边缘概率；Viterbi 用取最大合并路径并保存回溯指针，得到的是

\[
\widehat y=\argmax_y p_\theta(y\given x)=\argmax_y s_\theta(x,y).
\]

两者答案不一样。把每个位置边缘概率最大的标签拼起来，未必是得分最高的那条路径，甚至可能违反转移约束——回到开头那个 O 接 I-ORG 的毛病。要输出一条合法序列就得解码整条路径，要报告不确定性则要用边缘概率，别拿其中一个冒充另一个。

\subsection{梯度是观测计数减去模型期望}

给一对训练样本 \((x,y)\)，条件对数似然为

\[
\ell(\theta)=\theta^{\trans}F(x,y)-\log Z_\theta(x),
\qquad
F(x,y)=\sum_t f_t(y_{t-1},y_t,x,t),
\]

求导得到

\[
\nabla_\theta\ell(\theta)
=F(x,y)-\E_{p_\theta(y'\given x)}\bigl[F(x,y')\bigr].
\]

第一项是标注路径上的特征计数，第二项是模型在所有路径上的期望计数。梯度为零意味着两者相等，训练做的事是把模型的期望统计量拧到与数据一致。第二项看似要遍历 \(K^{T}\) 条路径，实际只需要前向--后向给出的相邻标签对边缘概率就能算出来。

这个``数据统计量减模型统计量''的形状与一般无向能量模型一致，差别在于 CRF 的配分函数可以精确算。稠密无向图上同一个式子的第二项算不动，只能退到采样或变分近似，这也是链结构值得珍惜的地方。加上 \(L_2\) 正则后，训练在抬高标注路径相对得分的同时，不让参数无限放大去追求单点确定性。

\subsection{全局归一化只修一类偏差}

局部归一化的序列模型在每个状态上单独算下一标签的概率分布。若某个状态的合法出口很少，它分配概率时几乎没得选，即使当前输入的证据很弱，也会把概率质量原样交出去；证据强的位置反而无法把之前的错误拉回来。这个现象叫标签偏置。CRF 对完整路径统一归一化，后面位置的证据可以重新参与整条路径的比较，这一类偏差随之消失。

它修不了别的。若真正的依赖跨越十几个词，而模型只有一阶转移和一堆手工特征，CRF 照样表达不出来；提高 Markov 阶数会让状态数成倍增长，\(K^2\) 变成 \(K^3\)、\(K^4\)。现在常见的分工是神经网络负责从上下文里算出每个位置的得分，CRF 层只管相邻标签的一致性。报告这类模型时值得把两边分开说清楚：网络编码了哪些上下文信息，CRF 约束了哪些输出结构。

\subsection{评估要回到结构化任务本身}

条件对数似然改善不保证下游指标改善。标签严重不平衡时，逐词准确率会被大量 O 标签撑起来——开头那个 \(96\%\) 就有相当一部分来自这里。逐词的置信度也不是整条路径正确的概率，后者要用路径概率或 \(n\)-best 的得分差来估。

该看的是实体级的精确率、召回率与 \(F_1\)，再单独检查非法转移是否清零、长实体是否被截断、未登录词上的表现，以及换一个来源的简历时性能掉多少。CRF 承认输出之间有结构，并在所有候选之间做联合比较，它不等于在分类器后面加一层后处理；只有当标签依赖确实重要、链上动态规划的代价又付得起时，它才比独立分类更值得用。

到这里，本单元已经出现了四种沿时间轴排列的模型：Markov 链、隐 Markov 模型、Kalman filter、条件随机场。它们的图长得很像，回答的却是不同的问题。下一节把这四者摆在一起，看清楚该按什么条件挑选。
